\uuid{ljh0}
\titre{Pgcd}
\auteur{bodin ; hanani ; mzari}
\organisation{exo7}
\licence{CC BY-NC-SA 4.0}
\datecreate{2020-01-20}
\niveau{L1}
\module{Algèbre}
\chapitre{Polynôme, fraction rationnelle}
\sousChapitre{Division euclidienne}
\difficulte{3}
\theme{Arithmétique des polynômes, pgcd}
\exo7qcmid{lille-1.polynomes.017}

\contenu{
\begin{question}
\qclassification{105.01,105.02}
\qtags{motcle=pgcd}
Soit $A(X) = X^4-X^2$. Soit $B(X) = X^2+X-2$.
Soit $D$ le pgcd de $A$ et $B$ dans $\Rr[X]$.
\begin{answers}
    \bad{$D(X) = 1$}

    \good{Il existe $U,V \in \Rr[X]$ tels que $AU+BV = X-1$.}

    \good{Il existe $u \in \Rr$ et $V \in \Rr[X]$ tels que $Au+BV = X-1$.}

    \bad{Il existe $U\in \Rr[X]$ et $v \in \Rr$ tels que $AU+Bv = X-1$.}
\end{answers}
\begin{explanations}
$A(X)=X^2(X-1)^2$, $B(X)=(X-1)(X+2)$, $D(X)=X-1$. $U(X)= -\frac14$, $V(X)=\frac14(X^2-X+2)$ donnent $AU+BV=D$.
\end{explanations}
\end{question}
}
